% === Begin Label Data ===
% ID: 520
% 难度星级: 
% 标签: 
% 备注: 
% 组卷引用次数: 0
% === End  Label Data ===

\begin{problem}{2021}{G}{新高考II卷}{4}{三角函数}
卫星导航系统中，地球静止同步轨道卫星的轨道位于地球赤道所在平面，轨道高度为 $ 36000\,\mathrm{km} $（轨道高度指卫星到地球表面的最短距离），把地球看成一个球心为 $ O $ 半径为 $ 64000\,\mathrm{km} $ 的球，其上点 $ A $ 的纬度是指 $ OA $ 与赤道所在平面所成角的度数，地球表面能直接观测到的一颗地球静止同步轨道卫星的点的纬度的最大值记为 $ \alpha $，该卫星信号覆盖的地球表面面积 $ S=2\pi r^{2}(1-\cos\alpha) $，（单位：$ \mathrm{km}^{2} $），则 $ S $ 占地球表面积的百分比为
\begin{choices}
\choice{{$26\%$}}
\choice{{$34\%$}}
\choice{{$42\%$}}
\choice{{$50\%$}}
\end{choices}
\end{problem}

\begin{answer}

\end{answer}

\begin{solutions}

\end{solutions}